\section{Co-registered Multilooked Complex (CMLC \& GMLC)}


\subsection{Product Description}

The multi-looked phase product \acs{cmlc} is going to be generated based on \acs{slc} products and can be later ground-projected and mosaicked for improved usability, or it may be directly generated from \acs{gslc} products.

Many of the characteristics of this product are shared with the INT and GINT products described above. All model-driven and data driven corrections (atmosphere, solid earth tides, etc.) can be applied to this product and are naturally inherited from the \acs{cslc} if the \acs{cslc} already includes them. 

The target posting is 50 m, according to the ROSE-L MRD. In the case of existing triangulated phase intermediate products, the posting might be in the order of 50-200~m depending on the sensor, whereas the number of averaged pixels can vary even within the same scene, for example when pixels are selected adaptively. More recent realizations make use of machine learning techniques to estimate the phase, so that it might not even be possible to determine the number of looks in a conventional way.

It is important for this product to include a quality layer which can be very useful for the user to guide phase unwrapping. Some quality layer is normally accompanying the triangulated phases we take inspiration from: for example \cite{Ferretti2011, Ansari2018} introduces coherence indices that describe succinctly the success of the triangulation procedure. Also in the procedure described in \cite{DeZan2022} there is a possibility to measure a pair of coherence indices. However, in this case, they vary for each point and each acquisition, so they do not describe synthetically the entire time series.

It is recommended to generate \acs{cmlc} products from \acs{cslc}s defined on a reference radar grid common to Sentinel-1 and ROSE-L missions for interoperability. In any case, the geocoded version (\acs{gmlc}) will be on the same grid for Sentinel-1 and ROSE-L.

Sub-band phases would be very useful to enable ionospheric corrections based on the split-spectrum technique. This means that we should consider generating 2-3 sets of phase time series for each dataset.

\subsubsection{Input and Auxiliary Data}
The input data for \acs{cmlc} is the \acs{cslc} product in radar geometry. The geocoded \acs{gmlc} product can be derived either by geocoding the \acs{cmlc} or directly from \acs{gslc} products. The former option allows more easily to work with a fixed number of looks which is considered preferable.

It would be advisable to incorporate ETAD-like corrections in the phase. This can happen implicitly when the corrections are already included in the \acs{cslc} input products.
If a more refined version of ETAD-like corrections is needed for high-end applications (e.g. precise strain measurement over large distances), this can be done at a later stage. To this end, it should be possible to undo the applied corrections.


\begin{longtblr}[
	label = tab:in_aux_mlc_products,
	entry = none,
	caption = {\acs{mlc} input and auxiliary data},
	]{
		width = \linewidth,
		colspec = {Q[l,20]Q[l,60]Q[c,20]},
		hlines,
		hline{1-Z} = {-}{0.05em},
	}
	\textbf{Name}   & \textbf{Description}   & \textbf{Mandatory (M) / Optional (O)} \\
	\acs{cslc} or \acs{gslc} & Co-registered (preferred) or geocoded L1 Single Look Complex image and annotations & M \\
    %\acs{dem} & External DEM (e.g., Copernicus DEM at different resolutions, with 30 m as the default worldwide). & M \\
    %Precise Orbits & \acs{slc} precise orbits & M \\
    \acs{etad} & \acs{etad}-like products for atmospheric delay and timing corrections & O \\
    \acs{tec} & External \acs{tec} data for delay correction, if not already included as an \acs{etad} layer & O \\
\end{longtblr}


\subsubsection{Main Variables}

The following tables (\ref{tab:var_cmlc_products} and \ref{tab:var_gmlc_products}) collect the main variables for the \acs{cmlc} and \acs{gmlc} products.

\begin{longtblr}[
	label = tab:var_cmlc_products,
	entry = none,
	caption = {\acs{cmlc} main variables},
	]{
		width = \linewidth,
		colspec = {Q[l,13]Q[l,42]Q[l,19]Q[l,14]Q[l,12]},
		hlines,
		hline{1-Z} = {-}{0.05em},
	}
	\textbf{Name} & \textbf{Description} & \textbf{Dimensions} & \textbf{Data Type} \\
	phase & Co-registered multi-look complex data.
       
    The phase corresponds conceptually to the interferometric phase between each acquisition and a fixed reference (e.g. the first acquisition in the stack) that one could compute directly by forming an interferogram with the reference. In practice, the \acs{mlc} phase should be less noisy. The \acs{mlc} phase can be used to generate directly multi-looked interferograms between any pair of acquisitions.
    & polarization, acquisition\_date, azimuth\_time, range\_time & complex64 (or float64) \\
    long-term coherence & The long-term coherence of each acquisition with a past reference. This information can be useful as an indication of the noisiness of interferograms computed between acquisitions that are separated by a large time interval (e.g. 1 year or more). The exact interpretation depends on the chosen processing parameters. & polarization, acquisition\_date, azimuth\_time, range\_time & float32 \\
    short-term coherence & The short-term coherence of each acquisition with the current reference (recent past). This information can be useful as an indication of the noisiness of interferograms computed between acquisitions that are separated by a short time interval (e.g. a few weeks), f.i. to guide phase unwrapping. The exact interpretation depends on the chosen processing parameters. & polarization, acquisition\_date, azimuth\_time, range\_time & float32 \\
\end{longtblr}

\begin{longtblr}[
	label = tab:var_gmlc_products,
	entry = none,
	caption = {\acs{gmlc} main variables},
	]{
		width = \linewidth,
		%colspec = {Q[l,15]Q[l,37]Q[l,20]Q[l,15]Q[l,13]},
        colspec = {Q[l,13]Q[l,42]Q[l,19]Q[l,14]Q[l,12]},
		hlines,
		hline{1-Z} = {-}{0.05em},
	}
	\textbf{Name} & \textbf{Description} & \textbf{Dimensions} & \textbf{Data Type} \\
	phase & Geocoded multi-look complex data:
       
    The phase corresponds conceptually to the interferometric phase between each acquisition and a fixed reference (e.g. the first acquisition in the stack) that one could compute directly by forming an interferogram with the reference. In practice, the \acs{mlc} phase should be less noisy. The \acs{mlc} phase can be used to generate directly multi-looked interferograms between any pair of acquisitions.
    & polarization, acquisition\_date, x, y & complex64 (or float64) \\
    long-term coherence & The long-term coherence of each acquisition with a past reference. This information can be useful as an indication of the noisiness of interferograms computed between acquisitions that are separated by a large time interval (e.g. 1 year or more). The exact interpretation depends on the chosen processing parameters. & polarization, acquisition\_date, x, y & float16 \\
    short-term coherence & The short-term coherence of each acquisition with the current reference (recent past). This information can be useful as an indication of the noisiness of interferograms computed between acquisitions that are separated by a short time interval (e.g. a few weeks), f.i. to guide phase unwrapping. The exact interpretation depends on the chosen processing parameters. & polarization, acquisition\_date, x, y & float16 \\
\end{longtblr}

The table below describes the dimensions/coordinates associated with the \acs{mlc} main variables.

\begin{longtblr}[
	label = tab:var_mlc_coordinates,
	entry = none,
	caption = {\acs{mlc} main variables coordinates/dimensions},
	]{
		width = \linewidth,
		colspec = {Q[l,20]Q[l,62]Q[l,18]},
		hlines,
		hline{1-Z} = {-}{0.05em},
	}
	\textbf{Name} & \textbf{Description} & \textbf{Data Type} \\
	polarization & Acquired polarizations & string \\
    band & Acquired band (e.g. main or auxiliary band) & string \\
    acquisition\_date & Image acquisition dates, formatted as YYYY-MM-DD & string \\
    azimuth\_time & Time along satellite track direction & datetime64 \\
    range\_time & Time across satellite track direction & float64 \\
    x & UTM/UPS x (East) coordinates [m] & float64 \\
    y & UTM/UPS y (North) coordinates [m] & float64 \\
\end{longtblr}


\subsubsection{Main Metadata}

The following table lists the main metadata and per-pixel metadata associated with the \acs{mlc} product. This product is not contemplated in \acs{ceos} documents nor in current SAR missions.  Many of the metadata are going to be identical for all the products in a stack and should not be repeated for each acquisition.

\begin{longtblr}[
	label = tab:meta_mlc_products,
	entry = none,
	caption = {\acs{cmlc}/\acs{gmlc} main metadata and per-pixel metadata},
	]{
		width = \linewidth,
		colspec = {Q[l,17]Q[l,38]Q[l,20]Q[l,13]Q[c,12]},
		hlines,
		hline{1-Z} = {-}{0.05em},
	}
	\textbf{Name} & \textbf{Description} & \textbf{Dimensions} & \textbf{Data Type} & \textbf{Per-pixel Metadata}\\
	input\_slc  & List of links to input reference \acs{slc} products & - & list of strings & \\
    reference\_dem & Link to reference DEM & - & string & \\
    phase\_demod & Phase to be applied to demodulate each image belonging to the \acs{slc} stack, computed using DEM & (x, y) or (azimuth\_time, range\_time) & float64 & \ding{51} \\ 
    Ionospheric phase & The ionospheric phase estimated with split-spectrum and subtracted from the phase. This variable allows undoing the correction for an advanced user who would want to utilize a different correction & (x, y) or (azimuth\_time, range\_time) & float32 & \ding{51} \\
    etad\_file & ETAD file used for each image belonging to the \acs{cslc}/\acs{slc} stack. This is optional: it is included only if used; otherwise, it is null & - & string & \\ 
    tec & Link to external \acs{tec} data, if used & - & string & \\ 
    look\_vector & 3-D components radar unit look vector, in an Earth-Centred Earth-Fixed (\acs{ecef}) coordinate system. It consists of unit vectors from antenna to surface pixel. Only for reference geometry. &  (x, y) or (azimuth\_time, range\_time) & float64 & \ding{51} \\
    inc\_angle & List of links to reference DEM-based local incidence angle maps. & - & string & \\

    
    %reference\_anf & Optional attribute: list of links to reference area normalization factor or local contributing area. & - & string & \\
    mask & Mask containing, for each ground cell, the following information:
    
     - 0: no distortions
    
     - 1: layover
    
     - 2: shadow
    
     - 3: undefined & (x, y) or (azimuth\_time, range\_time) & byte & \ding{51} \\
    nlooks & Number of equivalent looks for interferometric phase computation (only for reference geometry) & (x, y) or (azimuth\_time, range\_time) & float32 & \ding{51} \\
    
    processing\_
    
    parameters  & List of processing parameters & - & list of strings & \\
    units & Units of the variables of measurements.  & - & list of strings & \\
    
\end{longtblr}



\subsubsection{Reference System Tiling and Slicing}
\label{sec:mlc_reference_system_tiling}
\autoref{tab:mlc_ref_system_slicing} reports the reference system, the slicing, and grid for "Multi-Look Complex" products.


\begin{longtblr}[
	label=tab:mlc_ref_system_slicing,
	entry = none,
	caption = {Reference system and Slicing for \acs{cmlc} and \acs{gmlc}},
	]{
		width=\linewidth,
		colspec = {Q[l,15]Q[l,40]Q[l,15]Q[l,30]},
		%vline{2-5} = {-}{},
		%vline{5-5} = {-}{},
		hlines,
		hline{1-Z} = {-}{0.05em},
		rowhead=2,
	}
	\textbf{Product} & \textbf{Reference System} & \textbf{Slicing} & \textbf{Spacing}\\
	\textbf{\acs{cslc}} & Radar: azimuth/slant-range & Burst & sub-grid of original \acs{slc} radar grid, approximately 50 m \\
	\textbf{\acs{gslc}} & Cartographic: UTM/UPS & Burst & 50 m \\   
	
\end{longtblr}

\subsection{Product Requirements}

\subsubsection{Accuracy}
There is no standard praxis around the definition of requirements for multi-looked phase products. The most natural way of defining requirements would be to derive them from user requirements at the application level. 

Requirements on the product accuracy should ideally be specified in terms of bias and variance on the estimated phase or mean velocity. This is not trivial as it depends on the target correlation characteristics and user requirements do not have this level of detail.

One possibility is to define requirements with reference to a standardized coherence scenario that includes decorrelation, closure phases, and long-term coherence. 

Another possibility, which we prefer here, is to define the performance requirements relative to a full-covariance matrix estimator such as EMI or relative to an ideal performance such as the one predicted by the Cram\'er-Rao bound.

We are going to ask that the phase bias relative to a reference full-covariance algorithm should be smaller than 1 mm / year (TBC), which we set looking at the deformation precision requirement. 

The standard deviation degradation with respect to a reference full-covariance algorithm should be smaller than 3 dB (TBC).

The generation of this product requires a precise co-registration of the \acs{cslc} input products. It is important that the co-registration offsets are annotated in the \acs{cslc} and that it is possible to distinguish the geometrical/modeled part from the data-driven part (estimated typically with local cross-correlation).


\begin{longtblr}[
  label = tab:mlc_accuracy,
  entry = none,
  caption = {Accuracy Requirements for \acs{mlc} products},
]{
  width = \linewidth,
  colspec = {Q[l,35] Q[l,55] Q[c,20]},
  hlines,
  hline{1-Z} = {-}{0.05em},
  rowhead = 1,
}
\textbf{Metric} & \textbf{Description} & \textbf{Requirement} \\
$v_{\mathrm{bias}}$ & Deformation velocity bias relative to reference full-covariance method& $\leq 1$ mm/year \\
$20 \log_{10}(\sigma_\varphi/\sigma_\mathrm{ref})$ &  Phase std degradation relative to reference full-covariance method & $\leq 3$ dB \\
\end{longtblr}


\subsubsection{Timeliness}
The timeliness requirements depend on the applications which are intended to be covered by this product. If we ignore emergency requirements, assuming they can be covered by nearest-pair interferograms, the \acs{mlc} product can have a more relaxed latency requirement of e.g. 1 day or 1 week. For precise tectonic measurements, it is worth including the best reanalysis tropospheric corrections which are currently available only with a latency of a few months. %If there is access to predicted troposphere layers from ECMWF, one should quantify the degradation with respect to reanalysis datasets to determine whether they are sufficient for applications.

One could therefore foresee two releases of the product, the second being only a correction of the phase variable - not a full reprocessing. A first release would be within one to two days after the acquisition, to be able to include quality orbit information.
A second release would be as soon as final numerical weather reanalysis data are available (currently, three months for ERA-5 reanalysis).

\autoref{fig:cmlc_from_cslc} shows the flowchart for generating the \acs{cmlc} and \acs{gmlc} products from \acs{cslc} input.

\begin{figure}[H]
		\centering
		\resizebox{\textwidth}{!}{
		\begin{tikzpicture}[node distance=0.5cm and 2cm]
			
			% Define the style for the nodes
			% \tikzstyle{node} = [circle, draw, fill=blue!20, text centered, minimum size=1.5cm]
			
			% Draw nodes
			\node[rectin] (init) {\acs{cslc}};
			\node[block, right=of init] (phase_est) {Phase estim.};
			\node[rectout, right=of phase_est] (cmlc) {\acs{cmlc}};
			\node[block, right=of cmlc] (refine) {Refinement};
			\node[block, below=of cmlc] (geocode) {Geocoding};
			\node[rectout, below=of geocode] (gmlc) {\acs{gmlc}};
			\node[rectout, right=of refine] (cmlc2) {\acs{cmlc}};
			\node[rectaux, below=of refine] (models) {External models};
			\node[block, right=of gmlc] (refine2) {Refinement};
			\node[rectout, right=of refine2] (gmlc2) {\acs{gmlc}};
			
			% Draw edges with arrows
			\draw[->, thick] (init) -- (phase_est);
			\draw[->, thick] (phase_est) -- (cmlc);
			\draw[->, thick] (cmlc) -- (refine);
			\draw[->, thick] (refine) -- (cmlc2);
			\draw[->, thick] (models) -- (refine);
			\draw[->, thick] (cmlc) -- (geocode);
			\draw[->, thick] (geocode) -- (gmlc);
			\draw[->, thick] (models) -- (refine2);
			\draw[->, thick] (gmlc) -- (refine2);
			\draw[->, thick] (refine2) -- (gmlc2);
			
			% Draw edges with arrows
			%\draw[->, thick] (dem) -| (coreg);
			%\draw[->, thick] (orbit) -- (coreg);
			\draw[dotted, thick] (gmlc2) -- (cmlc2);
						
			% Timeline with linked markers and adaptive length
			\def\timelineStart{init}
			\def\timelineEnd{cmlc2}
			\draw[thick, ->] (\timelineStart.west |- 0,-6) -- (\timelineEnd.east |- 0,-6) node[right] {Time};
			\foreach \node/\label in {init/Start+(1-12h), phase_est/{<1h}, refine2/20-90d, gmlc2/<1h} {
				\draw (\node.south |- 0,-6) -- (\node.south |- 0,-6.5) node[below] {\label};
				\draw[dotted, thick] (\node.south) -- (\node.south |- 0, -6.25); % Connects the node to the timeline
			}
			
		\end{tikzpicture}
		}
		\caption{Sample timeline for the generation of \acs{cmlc} and \acs{gmlc} starting from \acs{cslc} products.}
		\label{fig:cmlc_from_cslc}
\end{figure}


\subsection {Usability, Integrability and Interoperability}

\autoref{tab:mlc_characteristics} summarizes usage characteristics of \acs{cmlc} and \acs{gmlc} products.
\begin{longtblr}[
  caption = {Usability, Integrability, and Interoperability of "Multi-looked Complex" Products},
  label = {tab:mlc_characteristics},
]{
  colspec = {Q[2cm] Q[12cm]},
  rowhead = 1,
  width = \linewidth
}
\hline
\textbf{Aspect} & \textbf{Description} \\
\hline
Usability &
\begin{itemize}
  \item Displacement Monitoring: useful for detecting large displacements such as landslides, earthquakes, mining subsidence, etc. with InSAR sensitivity. Usability to monitor ice motion needs to be demonstrated.
  \item Time series: naturally suited for time-series analysis of surface displacement.
  \item Advanced phase estimation: it is robust to deformation bias so that the user does not have to be a  specialist in phase estimation.
  \item Global to Local Monitoring: Suitable for regional or global-scale observations (e.g., tectonics) and local events (e.g., fast landslides).
  \item The \acs{mlc} wrapped phase can even be fed directly to a machine learning model trained for motion detection or parameter inversion~\cite{Biggs2022} or, possibly, phase unwrapping.
  \item Broad application domain: Useful in natural hazard analysis, volcanology, earthquake monitoring, urban subsidence, hydrology, etc.
\end{itemize}
\\
Integrability &
\begin{itemize}
  \item GIS Compatibility: Easily used within \acs{gis} platforms, due to geographic referencing and data format.
  \item Multi-Sensor Fusion: Integrates well with optical or other SAR projected products for land cover, surface change, or hazard mapping.
  \item Multi-Sensor compatibility: Can be derived from various sensors and integrated in joint analyses.
  \item Multi-sensor phase unwrapping: the \acs{gmlc} product can be used in a joint search for a phase-unwrapping solution, which can benefit from wavelength diversity.
\end{itemize}
\\
Interoperability &
\begin{itemize}
  \item Sensor-Independent Formats: Phase values projected geometry can be standardized across sensors for interoperability in displacement workflows.
  \item Geolocation Harmonization: When geocoded, products align spatially with other \acs{eo} datasets (e.g., NRB, optical, DEMs) to support multi-source analysis.
  \item Software Ecosystem Support: Usable in major \acs{eo} software platforms, supporting scientific and operational pipelines.
\end{itemize}
\\
\hline
\end{longtblr}

